\documentclass[11pt,twocolumn]{article}
\usepackage[utf8]{inputenc}
\usepackage[T1]{fontenc}
\usepackage{mathpazo}
\usepackage[margin=1.8cm,top=2cm,bottom=2cm]{geometry}
\usepackage{graphicx}
\usepackage{amsmath,amssymb}
\usepackage{booktabs}
\usepackage{xcolor}
\usepackage{hyperref}
\usepackage{enumitem}
\usepackage{titlesec}
\usepackage{fancyhdr}
\usepackage{wrapfig}
\usepackage{microtype}

\definecolor{quantablue}{RGB}{0,70,127}
\definecolor{accentgold}{RGB}{180,130,20}

\hypersetup{colorlinks=true,linkcolor=quantablue,urlcolor=quantablue,citecolor=quantablue}

\titleformat{\section}{\large\bfseries\color{quantablue}}{\thesection}{1em}{}
\titleformat{\subsection}{\normalsize\bfseries\color{quantablue!80!black}}{\thesubsection}{1em}{}

\pagestyle{fancy}
\fancyhf{}
\fancyhead[L]{\small\color{gray}\textit{Connectome Retrieval with Graph Neural Networks}}
\fancyhead[R]{\small\color{gray}\thepage}
\renewcommand{\headrulewidth}{0.4pt}

\title{\vspace{-1cm}\textbf{\color{quantablue}Reading Minds Backward:\\[4pt] How Graph Neural Networks Are Learning\\[2pt] to Reverse-Engineer the Brain's Wiring}}
\author{%
\textit{A perspective on the research program of Allier, Heinrich, Lappalainen, Schneider \& Saalfeld}\\[2pt]
\small Janelia Research Campus (HHMI) $\cdot$ Allen Institute $\cdot$ University of T\"ubingen\\[6pt]
\small\color{gray} Written in the style of \textit{Quanta Magazine} $\cdot$ March 2026%
}
\date{}

\begin{document}
\maketitle
\thispagestyle{fancy}

\begin{abstract}
\noindent
\textbf{A research program spanning three years and four biological organisms has demonstrated that graph neural networks can reverse-engineer neural circuit connectivity from activity recordings alone.} Starting from synthetic particle systems and advancing through simulated neural assemblies to biologically realistic connectomes of the fly visual system, the work reveals a counterintuitive principle: noise \emph{helps} solve the inverse problem. The approach recovers not just who connects to whom, but synaptic strengths, time constants, neuron types, activation functions, and hidden stimuli---achieving $R^2 > 0.99$ on circuits with 13,741 neurons and 434,122 synapses. A mathematically rigorous degeneracy analysis shows precisely when and why the problem is solvable, while an LLM-driven experimental loop automates the scientific discovery process itself.
\end{abstract}

\section{The Inverse Problem of the Brain}

Neuroscience has a data problem---but not the kind you might expect. Thanks to advances in electron microscopy and connectomics, we can now map the complete wiring diagram of entire nervous systems: every neuron, every synapse, every connection. The fruit fly \textit{Drosophila melanogaster} had its full brain connectome published in 2024---roughly 140,000 neurons and 50 million synaptic connections, painstakingly traced from nanometer-resolution imagery.

But knowing \emph{who} connects to \emph{whom} is only half the story. A wiring diagram is like having the street map of a city without knowing the traffic patterns. To understand how a neural circuit actually computes---how it transforms sensory inputs into behavior---you need to know the \emph{effective synaptic weights} (how strongly each connection operates), the \emph{time constants} (how quickly each neuron responds), and the \emph{nonlinear activation functions} (how signals are transformed at each synapse).

Can you run the problem backward? Given recordings of neural activity---the ``traffic''---can you infer the quantitative circuit model that generated it?

This is the question that Cedric Allier and collaborators at Janelia Research Campus have been systematically attacking since 2024. Their answer, refined across two published papers and an ambitious ongoing project spanning four biological organisms, is: \emph{yes, under precise mathematical conditions that we can characterize exactly.}

\section{From Particles to Neurons}

The story begins not with brains but with particles. In a 2024 paper~[1], the team demonstrated that \textbf{message-passing graph neural networks} (GNNs) could decompose the dynamics of heterogeneous interacting systems into interpretable components. The key insight was architectural: a GNN's computation---each node aggregating weighted messages from its neighbors---mirrors how physical systems actually work. Particles feel forces from neighbors. Neurons integrate synaptic inputs.

The framework was tested on seven systems: attraction-repulsion particles, gravitational dynamics, Coulomb interactions, flocking ``boids,'' wave propagation, reaction-diffusion patterns, and a signaling network. In every case, the GNN simultaneously recovered the interaction rules, the heterogeneous node properties (masses, charges, neuron types), and even the governing equations---the last step accomplished by feeding learned functions into symbolic regression, which returned $F \propto m/d^2$ and $F \propto q_1 q_2 / d^2$ without being told Newton's or Coulomb's laws.

The signaling network experiment---998 nodes, 17,865 edges, two node types---was the bridge. It showed that a GNN trained to predict one-step derivatives of node activity could recover the full connectivity matrix ($R^2 = 1.00$) and classify node types with perfect accuracy.

\section{Scaling to the Fly Brain}

The 2026 follow-up~[2] took this framework into neuroscience proper. The team designed a teacher-student paradigm: a biologically realistic ODE simulator (the ``teacher'') generates neural activity from known circuit parameters, and a GNN (the ``student'') learns to predict that activity while its learned edge weights are compared against ground truth.

The system studied is \emph{not} a toy model. The Drosophila visual system simulation, built on the \texttt{flyvis} model of Lappalainen et al.\ (2024), comprises \textbf{13,741 neurons} across \textbf{65 cell types} with \textbf{434,122 synaptic connections}, organized into 217 retinotopic columns. Each neuron follows a non-spiking ODE:
\begin{equation}
\tau_i \frac{dv_i}{dt} = -v_i + V_{\text{rest},i} + \sum_j W_{ij} \cdot \text{ReLU}(v_j) + I_i(t) + \sigma\eta_i
\label{eq:ode}
\end{equation}
where $\tau_i$ is the time constant, $V_{\text{rest},i}$ the resting potential, $W_{ij}$ the synaptic weight from neuron $j$ to $i$, $I_i(t)$ the visual stimulus, and $\sigma\eta_i$ process noise.

The GNN mirrors this structure with learnable components:
\begin{equation}
\frac{d\hat{v}_i}{dt} = f_\theta\!\left(v_i, \mathbf{a}_i,\; \sum_j \hat{W}_{ij} \cdot g_\phi(v_j, \mathbf{a}_j)^2,\; I_i\right)
\end{equation}
where $\hat{W}_{ij}$ are learnable edge weights, $g_\phi$ and $f_\theta$ are MLPs (the ``transfer'' and ``update'' functions), and $\mathbf{a}_i \in \mathbb{R}^2$ are learnable per-neuron embeddings that capture cell-type identity without supervision.

\section{The Noise Paradox}

The most striking finding is counterintuitive: \textbf{adding noise to the neural dynamics dramatically improves connectivity recovery.}

At zero noise, the GNN achieves $R^2_W = 0.90$ for synaptic weights---impressive, but imperfect. At moderate noise ($\sigma = 0.05$), recovery jumps to $R^2_W = 0.96$. At high noise ($\sigma = 0.5$), it reaches $R^2_W = 0.997$---near-perfect recovery of all 434,122 synaptic weights, along with time constants ($R^2_\tau = 1.00$), resting potentials ($R^2_{V_\text{rest}} = 0.96$), and cell-type clustering accuracy of 0.91 across 65 types.

Why does noise help? The answer lies in a rigorous degeneracy analysis. The fly's visual system is organized into retinotopic columns: the same 65 cell types are repeated across 217 spatial positions. In the noise-free regime, same-type neurons in different columns have nearly identical activation patterns. This creates a massive null space---approximately \textbf{121,100 of the 434,112 edge weights} (28\%) can be changed freely without affecting the observed dynamics \emph{at all}. The inverse problem is not merely difficult; it is mathematically ill-posed.

Process noise breaks this degeneracy. Each neuron receives independent stochastic perturbations, causing same-type neurons to diverge in their moment-to-moment activity. This lifts the rank of the activity matrix, collapsing the null space and making individual synaptic weights identifiable.

\section{Four Organisms, One Framework}

The current phase of the research---the \texttt{connectome-gnn} project---extends the approach to four biologically grounded neural circuits:

\begin{table}[h!]
\centering
\small
\begin{tabular}{@{}lrrp{2.8cm}@{}}
\toprule
\textbf{Circuit} & \textbf{Neurons} & \textbf{Edges} & \textbf{Dynamics} \\
\midrule
Fly visual system & 13,741 & 434K & ReLU, graded \\
Fly central complex & 152 & 9,722 & Softplus ring attractor \\
Larva motor circuit & 230 & 4,222 & Two-population, gain-mod. \\
Zebrafish oculomotor & 609 & 10,665 & Pure linear integrator \\
\bottomrule
\end{tabular}
\end{table}

Each circuit presents distinct challenges. The central complex is a ring attractor with softplus nonlinearity and per-neuron gains. The larva has a two-population architecture (premotor and motor neurons). The zebrafish is perhaps the hardest: its purely linear dynamics mean that infinitely many connectivity matrices produce identical activity, making the inverse problem fundamentally degenerate without noise.

Results across all four systems confirm the noise-helps principle. The zebrafish is the most dramatic example: connectivity $R^2$ rises from 0.02 (noise-free, fully connected) to 0.99 (noise = 0.5).

\section{What You Can and Cannot Recover}

A Procrustes-aligned SVD analysis of the learned connectivity matrices reveals a telling asymmetry. Decomposing $W = U \Sigma V^T$, the \emph{output connectivity modes} $U$ (which neurons co-activate) are consistently well-recovered ($R^2 \approx 0.90$ for the central complex), while the \emph{input selection modes} $V$ (which neurons drive which) are much harder ($R^2 \approx 0.51$).

This matches a theoretical prediction by Mastrogiuseppe \& Ostojic (2018): in recurrent networks, the output patterns are directly constrained by observed dynamics, while input selection is only indirectly observable. The $U$--$V$ gap widens precisely when overall recovery is harder, providing a diagnostic for identifiability.

The method is also remarkably robust to corrupted prior knowledge. When the adjacency matrix is contaminated with 100--400\% extra false edges, the GNN learns to assign near-zero weights to spurious connections, maintaining $R^2_W > 0.94$. Even with 20\% of true edges missing, recovery remains viable ($R^2_W = 0.80$).

\section{The LLM Scientist}

Perhaps the most unconventional aspect of this research is how experiments are designed. The team built \texttt{GNN\_LLM.py}, an automated pipeline where an LLM (Claude) acts as an iterative experimental scientist:

\begin{enumerate}[leftmargin=*,itemsep=2pt]
\item The LLM reads previous experimental results and accumulated principles.
\item It proposes a new hyperparameter configuration, reasoning about which parameters to change and why.
\item The experiment runs (data generation, GNN training, evaluation).
\item Results are fed back; the LLM updates its working memory.
\end{enumerate}

Over 35 exploration campaigns (128 iterations each) across all four biomodels, this agentic loop discovered model-specific configurations that significantly outperform hand-tuned defaults: $+22\%$ $R^2_W$ for noise-free flyvis, $+38\%$ at moderate noise. It identified that Dale's law regularization is essential for the central complex but destructive for the larva, that learning rate on $W$ and L1 sparsity are the dominant levers, and that many architectural parameters (MLP depth, hidden dimension) have surprisingly flat dose-response curves.

The LLM doesn't just tune hyperparameters---it accumulates \emph{causal understanding} of why certain regimes succeed or fail, discovering phenomena like bimodal convergence landscapes and sharp augmentation cliffs that would be invisible to grid search.

\section{Beyond Simulation: The Road Ahead}

The ultimate goal is real neural data: calcium imaging, electrophysiology, light-sheet recordings. The team has already begun experiments with the Zapbench zebrafish electrophysiology dataset and is preparing a NeurIPS 2025 submission targeting the full four-organism benchmark with baseline comparisons (MLP, RNN, Neural ODE, state-space models).

Critical open questions remain. Can the framework handle time-dependent connectivity (synaptic plasticity)? Time delays? The gap between voltage traces and calcium indicators? Missing neurons? The degeneracy analysis provides a precise mathematical framework for answering these questions: each additional biological complication either expands or contracts the null space, and the answer depends on the specific circuit architecture.

\vspace{0.4cm}
\noindent\rule{\columnwidth}{0.6pt}

\noindent\textbf{References}

\noindent{\small [1] Allier, C., Schneider, M.C., Innerberger, M., Heinrich, L., Bogovic, J.A., \& Saalfeld, S. (2024). Decomposing heterogeneous dynamical systems with graph neural networks. \textit{arXiv:2407.19160}.

\noindent [2] Allier, C., Heinrich, L., Schneider, M., \& Saalfeld, S. (2026). Graph neural networks uncover structure and functions underlying the activity of simulated neural assemblies. \textit{arXiv:2602.13325}.

\noindent [3] Lappalainen, J.K., et al.\ (2024). Connectome-constrained networks predict neural activity across the fly visual system. \textit{Nature}, 634, 1132--1140.

\noindent [4] Beiran, M.\ \& Litwin-Kumar, A.\ (2023). Parametric control of flexible timing through low-dimensional neural manifolds. \textit{Neuron}, 111, 739--753.}

\vspace{0.6cm}
\noindent\rule{\columnwidth}{0.6pt}

\section*{\color{quantablue}Perspective for ML Journal Submission}

\vspace{0.2cm}

\noindent\textit{The following paragraph is written in the style of a forward-looking ML journal perspective piece.}

\vspace{0.3cm}

\noindent Current graph neural network approaches to the neural inverse problem---recovering effective connectivity from activity data---will eventually be recognized as early steps toward a much more sophisticated future. That future will involve new activity-based and structural ingredients beyond pairwise synaptic weights: dendritic compartment models, gap junctions, neuromodulatory volumes, and time-varying plasticity rules. More generalized inference schemes will emerge, based on a variety of reference dynamics---not just teacher-student simulations but transfer from partially observed \emph{in vivo} recordings, multi-modal data fusion (calcium imaging, electrophysiology, connectomics), and hybrid mechanistic-learned models that interpolate between known biophysics and data-driven flexibility. The energy (or loss) landscape will increasingly incorporate machine-learned priors on circuit motifs, Dale's law, spectral constraints, and low-rank structure, guided by automated experimental design loops where LLM agents propose, execute, and reason about experiments at a pace no human team can match.

The rate-determining step in this transition will be the creation and curation of larger, higher-quality, and more comprehensive benchmarks at the interface of connectomics and physiology. Such benchmarks will need to span multiple organisms (from \textit{C.~elegans} and \textit{Drosophila} larvae to zebrafish and mouse cortex), multiple recording modalities (voltage imaging, calcium indicators, multi-electrode arrays), and multiple dynamical regimes (spontaneous activity, stimulus-driven, behavioral). They must be systematically designed to probe identifiability boundaries---the precise conditions under which the inverse problem transitions from well-posed to degenerate---well documented with ground-truth connectivity from electron microscopy, and readily accessible to the community as standardized, versioned datasets with associated evaluation protocols. Only with such infrastructure can the field move from proving that connectome recovery \emph{is possible in principle} to establishing \emph{when, why, and how reliably} it works in practice---and ultimately close the gap between the structural connectome and the functional effectome that governs neural computation.

\end{document}
